\section{Conclusion}
\label{sec:conclusion}

This paper introduced Hypothetica, a multi-agent system for automated originality assessment of research proposals. By integrating Large Language Models with Retrieval Augmented Generation across three heterogeneous corpora -- academic literature (OpenAlex), industrial prior art (Google Patents), and open-source repositories (GitHub) -- the system performs cross-domain semantic similarity analysis and produces interpretable, criterion-level originality estimates.

Benchmark evaluation across 150 manually labeled cases demonstrates a mean classification accuracy of 76.0\%, representing a $2.3\times$ improvement over random baseline. Performance is consistent across corpus types, with per-class accuracy particularly strong for the \textit{already\_exists} and \textit{novel} labels. The incremental class remains the most challenging across all corpora, reflecting both the inherent ambiguity of the boundary and the deliberate inclusion of borderline cases in the benchmark design. Critically, no catastrophic misclassifications were observed: all errors occurred between adjacent classes, preserving the ordinality of the originality spectrum.

The evaluation also surfaces a fundamental tension in novelty benchmark construction: realistic novel ideas necessarily share vocabulary, problem framing, and domain context with existing work. A system that penalizes this partial overlap -- rather than ignoring it -- provides more honest and practically useful feedback than one optimized purely for recall on the novel label. This positions Hypothetica not as a pass/fail filter but as a diagnostic tool that surfaces the specific dimensions of overlap and novelty in a research proposal, supporting researchers in iterating toward more clearly differentiated contributions.

Future work will focus on three directions. First, resolving the Likert score compression behavior in Layer~1 through rubric refinement and calibrated prompting, which we identify as the primary path toward breaking the current 76--78\% accuracy ceiling. Second, expanding retrieval coverage to include ACM Digital Library and IEEE Xplore, improving domain specificity for computer science and engineering literature. Third, developing longitudinal tracking capabilities to monitor how the novelty landscape of a research area evolves over time, enabling the system to flag ideas that were novel at submission but have since been overtaken by the literature.

Ultimately, Hypothetica contributes a transparent, evidence-grounded approach to originality evaluation -- one that bridges the gap between theoretical novelty and the rigorous prior art search that peer review demands, while remaining practically deployable at the scale of individual researchers and institutional screening workflows.
